% ------------------------------------------------------------------------
%
\begin{frame}{Proving associativity of stack concatenation}

The definition of the concatenation of two stacks was:
\begin{equation*}
\fun{cat}(        \el,t) \xrightarrow{\smash{\alpha}} t;\qquad
\fun{cat}(\cons{x}{s},t) \xrightarrow{\smash{\beta}}
\cons{x}{\fun{cat}(s,t)}.
\end{equation*}
Let us prove the associativity of \fun{cat}, which we express
formally as
\begin{equation*}
  \pred{CatAssoc}{s,t,u} \colon
\fun{cat}(s,\fun{cat}(t,u)) \equiv
\fun{cat}(\fun{cat}(s,t),u)
\end{equation*}
for all stack values \(s\), \(t\) and~\(u\).

\end{frame}

% ------------------------------------------------------------------------
%
\begin{frame}{Proving associativity of stack concatenation}

  We apply the well\hyp{}founded induction principle to the structure
  of~\(s\), so we must establish
  \begin{itemize}

    \item the basis \(\forall t,u \in S.\pred{CatAssoc}{\el,t,u}\),
      and the inductive

    \item step \(\forall s,t,u \in S.\pred{CatAssoc}{s,t,u}
      \Rightarrow \forall x \in T.\pred{CatAssoc}{\cons{x}{s},t,u}\).

  \end{itemize}

  \bigskip

  Underlining the call being rewritten, the base case is
  straightforward:
  \begin{equation*}
    \ufun{cat}(\el,\fun{cat}(t,u))
    \xrightarrow{\smash{\alpha}} \fun{cat}(t,u)
    \xleftarrow{\smash{\alpha}} \fun{cat}(\ufun{cat}(\el,t),u).
  \end{equation*}

  Let us assume now \(\pred{CatAssoc}{s,t,u}\), called the
  \textbf{induction hypothesis}, and let us prove
  \(\pred{CatAssoc}{\cons{x}{s},t,u}\), for any term~\(x\).

\end{frame}

% ------------------------------------------------------------------------
%
\begin{frame}{Proving associativity of stack concatenation}

  The goal here is to use the rewrite system as an abstract machine to
  prove a property, with the timely help of the induction principle.

  \bigskip

  Precisely, we want to rewrite each side of the equivalence we wish
  to prove until we either find the same term (equality), or we use
  the induction hypothesis (equivalence).

  \bigskip

  Since we are going to work with equivalences, we should lift the
  call\hyp{}by\hyp{}value reduction strategy here, as it does not
  matter when arguments are evaluated, as long as their evaluation
  terminates and is confluent: we always assume these properties on
  the system hold when proving properties of the functions it defines.

\end{frame}

% ------------------------------------------------------------------------
%
\begin{frame}{Proving associativity of stack concatenation}

  We may start with the left\hyp{}hand side of
  \(\pred{CatAssoc}{\cons{x}{s},t,u}\), in other words:
  \(\fun{cat}(\cons{x}{s},\fun{cat}(t,u))\).

  \bigskip

  We notice that this is \textbf{not} a value, as the call contains a
  reducible call, \(\fun{cat}(t,u)\).

  \bigskip

  When using functions to compute, we match a call containing only
  values against the patterns of the rules.

  \bigskip

  Here, we would like to match a pattern against a pattern. This is
  more general than matching, because values are patterns, and is
  called \textbf{unification}. Given two patterns, we want to find
  whether we can find assignments to the variables of these so that,
  when they are substituted by their denotation, the two patterns
  become equal.

\end{frame}

% ------------------------------------------------------------------------
%
\begin{frame}{Proving associativity of stack concatenation}

  Here, we would like to unify the \textbf{goal}
  \(\fun{cat}(\cons{x}{s},\fun{cat}(t,u))\) with the \textbf{pattern}
  \(\fun{cat}(\cons{x}{s},t)\).

  \bigskip

  The problem is that we have variables that are the same in the goal
  and the pattern (\(x\), \(s\)~and~\(t\)).

  \bigskip

  The case of \(\cons{x}{s}\) and \(\cons{x}{s}\) is harmless, but the
  unification of \(\fun{cat}(t,u)\) and~\(t\) is not.

  \bigskip

  That is why it is best to rename the variables in the pattern so
  they are not found in the goal (this is the
  \textbf{\(\alpha\)\hyp{}renaming} of \(\lambda\)\hyp{}calculus), and
  then try to unify them.

  \bigskip

  To simplify things even more, we are going to rename \emph{all} the
  variables in the definition of \fun{cat}.

\end{frame}

% ------------------------------------------------------------------------
%
\begin{frame}{Proving associativity of stack concatenation}

  Let us try for example the following renaming:
  \begin{align*}
    \fun{cat}(        \el,b) &\xrightarrow{\alpha} b\\
    \fun{cat}(\cons{i}{a},b) &\xrightarrow{\smash{\beta}}
    \cons{i}{\fun{cat}(a,b)}
  \end{align*}

  We want to unify the goal \(\fun{cat}(\cons{x}{s},\fun{cat}(t,u))\)
  with the pattern \(\fun{cat}(\cons{i}{a},b)\).

  \bigskip

  The unification yields the equations \(x = i\), \(s = a\) and
  \(\fun{cat}(t,u) = b\).

  \bigskip

  Then we apply the set of equations from the unification, so the
  right\hyp{}hand side of~\(\beta\), \(\cons{i}{\fun{cat}(a,b)}\),
  becomes \(\cons{x}{\fun{cat}(s,\fun{cat}(t,u))}\).

\end{frame}

% ------------------------------------------------------------------------
%
\begin{frame}{Proving associativity of stack concatenation}

  We have:
  \begin{equation*}
    \begin{array}{r@{\;}l@{\;}l@{\qquad}r@{}}
      \ufun{cat}(\cons{x}{s},\fun{cat}(t,u))
      & \xrightarrow{\smash\beta}
      & \cons{x}{\fun{cat}(s,\fun{cat}(t,u))}\\
      & \equiv
      & \cons{x}{\fun{cat}(\fun{cat}(s,t),u)}
      & \IH{\pred{CatAssoc}{s,t,u}}\\
      & \xleftarrow{\smash\beta}
      & \ufun{cat}(\cons{x}{\fun{cat}(s,t)},u)\\
      & \xleftarrow{\smash\beta}
      & \fun{cat}(\ufun{cat}(\cons{x}{s},t),u).
    \end{array}
  \end{equation*}

  Note the use of the inductive hypothesis \(\pred{CatAssoc}{s,t,u}\).

  \bigskip

  This derivation above establishes that
  \(\pred{CatAssoc}{\cons{x}{s},t,u}\) holds.

  \bigskip

  The induction principle entails that \(\forall s,t,u \in
  S.\pred{CatAssoc}{s,t,u}\).\hfill\(\Box\)

\end{frame}
